% Ce chapitre repose sur :
%   DOC: s30
%   OBS: presentation-du-besoin, rattachement-du-travail, organisation-en-phases
%   HYP: (aucun)
%   SECTIONS: 6

\chapter{Cadrage et méthodologie}\label{chap:cadrage}

\section{Le besoin exprimé}

Le travail répond à un cahier de charges portant sur l'activité du service d'accueil et
d'admission dans son ensemble, et non sur le seul processus observable depuis un poste. Ce
périmètre est déterminant pour la suite : il oblige à documenter quatre processus auxquels
l'observation directe n'a pas donné accès, et il explique la place que tient dans ce projet la
reconstruction par voie documentaire.

Le document lui-même ne figure pas au dépôt et n'a pas vocation à y figurer. Ce qui y figure, en
revanche, ce sont trois de ses exigences, consignées à l'endroit où elles ont été mises en défaut :
tout indicateur affiché doit être recalculé depuis les tables de faits, le tableau de bord ne doit
lire que la base, et les valeurs affichées doivent être lisibles par un lecteur qui ne connaît pas
les codes du système d'origine.

Les trois sont tenues en partie et manquées en partie, et les écarts sont écrits plutôt que
contournés. Six indicateurs ne sont pas recalculés depuis les faits, dont cinq parce qu'aucune table
de faits ne porte la matière — il n'existe ni fait des créances ni fait des relances — et un parce
que le recalcul changerait sa valeur en changeant sa population. La capacité litière, dont dépendent
quatre indicateurs de séjour, n'est lisible nulle part dans la base et transite par une table de
paramètres portant sa provenance. Enfin quatre dimensions sur six affichent leur code nu, faute de
libellé qu'une source établisse — et l'abstention est délibérée : un code nu est illisible et le
montre, un libellé inventé est lisible et le cache.

\aRediger{presentation-du-besoin}{la façon dont le besoin a été présenté au début du stage, et ce
que l'encadrement a précisé oralement sur les attentes et les priorités}\releve{presentation-du-besoin}

\section{Périmètre d'observation}

Le système d'information cloisonne l'accès par profil applicatif, et le service n'accorde à un
stagiaire qu'un seul des cinq profils que porte l'écran de démarrage. \textbf{Un processus a donc été
observé, quatre ont été reconstruits par voie documentaire.} Ce n'est pas une commodité : c'est la
contrainte à partir de laquelle tout le dispositif de provenance a été bâti.

Une observation et une reconstruction n'ont pas la même valeur, et les confondre reviendrait à
présenter comme vu ce qui a été déduit. Le projet ne les confond pas : \textbf{chaque colonne du modèle de
données porte une étiquette de provenance, et l'exigence de preuve diffère par étiquette.}

\begin{center}
\begin{tabular}{@{}p{1.6cm}p{5.0cm}p{1.6cm}p{5.5cm}@{}}
\toprule
Étiquette & Sens & Entrées & Ce que la preuve doit être \\
\midrule
\texttt{OBS} & champ observé à l'écran & \textbf{81} & une référence de relevé, et le libellé reproduit à l'identique \\
\addlinespace
\texttt{DOC} & champ déduit d'une source & \textbf{72} & un identifiant du registre des sources \\
\addlinespace
\texttt{HYP} & champ posé sans preuve externe & \textbf{22} & la mention explicite qu'aucune preuve n'existe \\
\bottomrule
\end{tabular}
\end{center}

\begin{center}
\footnotesize Les trois étiquettes de provenance et leur exigence de preuve, sur les 175 colonnes de
la couche source.
\end{center}

\medskip

Les vingt-deux colonnes hypothétiques portent toutes la mention d'absence de preuve externe, et ce
décompte est mesuré, non supposé. La proportion elle-même est un résultat : \textbf{un peu moins de la
moitié des colonnes est observée}, et le lecteur sait laquelle.

Le même dispositif s'étend à la prose du présent rapport. Chaque affirmation y relève d'une des trois
mêmes catégories — une source documentaire, un relevé d'observation, une convention posée —, chaque
fichier de chapitre déclare en tête ce sur quoi il repose, et un contrôle vérifie que la déclaration
coïncide exactement avec ce que le fichier porte, dans les deux sens.

\textbf{Aucune image du système d'information n'entre dans ce rapport, et le motif tient en une
phrase :} les écrans d'un dossier patient en production affichent des identifiants de session, des
libellés de compte et la structure des dossiers d'un établissement réel, et le dépôt de ce projet est
public. Ce que les images auraient montré, le relevé de champs le porte sous une forme plus
exploitable — quatre écrans, cent quarante-trois éléments, chacun avec son identifiant, son libellé
exact, son type apparent et sa valeur observée le cas échéant. Deux schémas d'écran redessinés
restituent l'agencement des blocs, sans aucun contenu. La règle est tenue par deux dispositifs
indépendants : une exclusion dans la configuration du gestionnaire de versions, et un contrôle
bloquant qui échoue si un fichier suivi porte une extension d'image sous le répertoire du rapport ou
à la racine — le premier seul ne protégerait pas d'une indexation forcée.

\section{La contrainte de confidentialité et la décision de générer}

Une décision de cadrage domine tout le projet : \textbf{aucune donnée réelle ne sort du service.} Elle
n'est pas négociable et elle n'a pas été négociée. Elle interdit d'extraire un échantillon, fût-il
anonymisé, et elle interdit donc de bâtir une chaîne de données sur des lignes venues du système
observé.

La conséquence est qu'un jeu de données est produit. Sa volumétrie n'est pas inventée : elle est
relevée par établissement nommé dans le recueil statistique national \cite{s30} et mise à l'échelle
de la période. Ses relations internes ne sont pas davantage improvisées — elles sont injectées
délibérément, une par une, et chacune est consignée avec son paramètre, son statut de provenance et
sa conséquence sur ce que le tableau de bord donne à voir.

\textbf{Ce que la génération permet.} Elle permet de démontrer une chaîne complète sur des volumes
réalistes, d'y injecter des défauts dont on connaît la vérité terrain, et donc de mesurer ce qu'un
traitement retrouve. Un rapprochement probabiliste évalué contre des doublons dont on sait
lesquels sont vrais est mesurable ; le même rapprochement sur des données réelles ne le serait pas.

\textbf{Ce qu'elle interdit de conclure, et c'est le point le plus important de cette section.} Aucune
grandeur produite par la chaîne n'est une mesure de l'établissement. Une relation injectée se
retrouve nécessairement à la sortie : la lire comme un résultat serait circulaire. Le projet écrit
cette limite là où elle se lit — le registre des relations injectées porte, pour chacune, ce que la
page affiche et ce qu'on n'a pas le droit d'en conclure. Le chapitre de recommandations en tire une
conséquence de forme, en séparant systématiquement ce qui établit un constat de ce qu'un indicateur
démontre.

\section{Démarche}

Deux principes gouvernent le travail, et l'un et l'autre se lisent dans le dépôt plutôt que dans une
déclaration d'intention.

\textbf{Mesurer avant de corriger.} Aucune correction n'est proposée avant que sa cause soit
mesurée, et aucun chiffre n'est écrit sans la commande qui le produit. Le principe se prolonge
au-delà de l'écriture : les enregistrements de décision portent, pour plusieurs d'entre eux, une
prémisse explicitement datée et la consigne de la re-mesurer avant toute modification qui la
mettrait en jeu. Une décision n'y est jamais présentée comme définitive, mais comme vraie sous des
conditions nommées.

\textbf{Vérifier par mutation.} Une propriété nouvellement affirmée n'est pas tenue pour vérifiée
parce qu'un contrôle est vert : le contrôle est délibérément mis en défaut, on s'assure qu'il
rougit, et l'état initial est restauré. Un contrôle vert qui reste vert quand on casse ce qu'il
surveille ne surveille rien.

La méthode a payé, et il faut dire ce qu'elle a trouvé plutôt que d'en vanter le principe. Elle a
révélé qu'un contrôle de correspondance passait par accident sur un fichier dépourvu de déclaration
— ensemble déclaré vide contre ensemble porté vide —, ce qui a fait ajouter l'exigence d'une
déclaration explicite. Elle a révélé ensuite que le retrait d'une section entière à l'intérieur d'un
chapitre n'était détecté que si cette section emportait une source qu'aucune autre ne portait :
mesuré sur huit sections, six étaient détectées et deux ne l'étaient pas, ce qui a fait déclarer le
décompte de sections. Elle a enfin établi qu'une vérification prévue ne pouvait pas mordre — le
convertisseur bibliographique ne lisant jamais le champ que la mutation visait —, cas où la
vérification est remplacée par celle qui a un sens plutôt que forcée.

\aRediger{rattachement-du-travail}{le rattachement du travail effectué aux phases du projet, et ce
que chaque phase a produit}\releve{rattachement-du-travail}

\section{Outils et justification}

Chaque brique de la pile technique est retenue pour un motif écrit, et non par habitude.

\begin{center}
\begin{tabular}{@{}p{3.4cm}p{4.0cm}p{6.4cm}@{}}
\toprule
Brique & Rôle & Motif du choix \\
\midrule
Serveur de base de données relationnelle & couche de persistance & un moteur client-serveur, non un moteur embarqué : trois services attaquent la même base par le réseau, avec des identifiants distincts, et le chargement écrit pendant que le tableau de bord lit \\
\addlinespace
Schéma en étoile & couche analytique & des faits entourés de dimensions, forme que l'interrogation analytique exploite et que le dossier doit démontrer \\
\addlinespace
Outil de transformation déclaratif & modélisation & matérialisation en vues, configuration de connexion hors dépôt, macros de conversion adossées aux formats mesurés à l'extraction \\
\addlinespace
Bibliothèque de rapprochement probabiliste & identités & un modèle de vraisemblance dont les paramètres s'estiment sur les données, plutôt qu'une collision exacte qui ne retrouve que l'identique \\
\addlinespace
Moteur en mémoire pour le rapprochement & exécution & le rapprochement écrit des tables intermédiaires nombreuses ; les tenir hors du serveur évite d'y laisser des résidus \\
\addlinespace
Orchestrateur de graphes de tâches & ordonnancement & le conteneur ordonnance sans exécuter : les tâches sont des commandes qui s'exécutent dans l'environnement du projet \\
\addlinespace
Bibliothèque d'application web analytique & tableau de bord & un service qui lit un schéma dédié, rafraîchi par échange de noms en une transaction \\
\bottomrule
\end{tabular}
\end{center}

\begin{center}
\footnotesize Pile technique et motif de chaque choix. Chacun est consigné dans un enregistrement
de décision distinct, avec la mesure qui l'a tranché et ce qui l'aurait invalidé.
\end{center}

\medskip

Deux traits de cette liste méritent d'être relevés. Le premier est que le motif d'un choix n'est
presque jamais la performance : c'est la propriété que le dossier doit démontrer, ou la contrainte
qu'il faut ne pas violer. Le second est que chaque enregistrement porte une rubrique disant ce qui
aurait invalidé la décision — un choix qu'on ne sait pas décrire comme réfutable n'est pas un choix,
c'est une préférence.

\section{Organisation du projet en phases}

\aRediger{organisation-en-phases}{les phases du projet, nommées par ce qu'elles font, avec pour
chacune son objectif et ses livrables, suivies d'une figure restituant leurs dépendances — aucun
calendrier, aucun relevé de jours}\releve{organisation-en-phases}
